\begin{abstract}
While post-hoc model merging has emerged as a powerful paradigm for combining task-specific expert models without retraining, adaptive test-time ensembling methods remain highly susceptible to severe transductive overfitting on local stream noise. In this paper, we establish the first rigorous statistical learning-theoretic foundation for adaptive model merging through our proposed framework, Rademacher-Bounded Polynomial Merging (RBPM). By constraining layer-wise merging trajectories to a low-degree polynomial subspace and introducing an analytical Rademacher complexity regularization penalty, RBPM provides provable out-of-distribution generalization guarantees. Our theoretical analysis demonstrates that the polynomial trajectory constraint reduces the empirical Rademacher complexity of the merging hypothesis space by a factor of $\mathcal{O}(\sqrt{L / \log(d)})$, where $L$ is the network depth and $d$ is the polynomial degree. Extensive empirical validation on a 12-layer deep CNN backbone across four distinct classification tasks confirms our theoretical findings. While unconstrained offline tuning achieves an average test accuracy of 32.75\%, RBPM successfully regularizes ensembling parameter capacity to achieve a robust average test accuracy of 38.85\%, representing a significant gain over the zero-optimization uniform ensembling baseline (29.05\%) while retaining strict learning-theoretic control over the generalization gap.
\end{abstract}